\clearpage
\section{Dataset Description}\label{sec:data}
\subsection{Source and task labels}
We use the Large Movie Review Dataset~\cite{maas2011}, distributed as
\code{stanfordnlp/imdb} on Hugging Face~\cite{imdbhub}. It provides 25,000 labeled
training reviews, 25,000 labeled test reviews and 50,000 unlabeled reviews. The
unlabeled portion is excluded from our supervised experiments. Labels are
\code{0 = negative} and \code{1 = positive}. The source benchmark defines
negative reviews from ratings of at most 4/10 and positive reviews from ratings
of at least 7/10; the exported records do not include a separate rating field.

The preparation notebook was executed on Google Colab CPU. The initial download
did not pin a dataset revision, so the verified exported files and their SHA-256
checksums define the data used by the team. This distinction matters when
reproducing the experiment from a changing remote dataset distribution.

\subsection{Quality audit and fixed partitions}
The original labeled splits contain no missing, empty or non-string text, and no
missing or invalid labels. Exact matching finds 96 extra duplicate rows within
training, 199 within test, and 123 distinct review texts shared by the original
training and test sets. No repeated text has conflicting labels.

We preserve the official test split. From training, we first remove the 123
overlapping records, then remove 96 remaining duplicate extra rows, keeping the
first occurrence. This leaves 24,781 eligible reviews. A label-stratified 80/20
split with seed 42 yields the fixed training and validation sets in
Table~\ref{tab:splits}. Test text is used for the exact-overlap audit, but not for
choosing preprocessing rules. Stable review IDs and a split manifest maintain
membership across all five pipelines.

\begin{table}[htbp]\centering
\caption{Fixed project partitions after training-side duplicate removal.}\label{tab:splits}
\begin{tabular}{lrrrr}\toprule
Split & Reviews & Negative & Positive & Negative (\%)\\\midrule
Training & 19,824 & 9,869 & 9,955 & 49.78\\
Validation & 4,957 & 2,468 & 2,489 & 49.79\\
Test & 25,000 & 12,500 & 12,500 & 50.00\\\midrule
Total & 49,781 & 24,837 & 24,944 & 49.89\\\bottomrule
\end{tabular}
\end{table}
\reportfigure{person1_image1.png}{0.88}{Label counts in the fixed project splits. Classes are approximately balanced; the full official test set is retained.}{fig:labels}

\subsection{Review lengths and examples}
Training reviews have a mean length of 234.70 whitespace-separated words, a median
of 175 and a standard deviation of 174.29. The minimum is 10, the maximum is
2,470, the interquartile range is 127--285, and the 95th percentile is 600.
These are word counts, distinct from spaCy tokens and transformer wordpieces.
The distribution is right-skewed (Figure~\ref{fig:lengths}). For readability,
199 reviews above the 99th-percentile cutoff of 917.77 words are omitted from
the plot only; they remain in training and in every summary statistic.
\reportfigure{person1_image2.png}{0.94}{Training review lengths before linguistic preprocessing. The chart omits the longest 1\% for display only.}{fig:lengths}

\begin{table}[htbp]\centering\small
\caption{Reproducibly sampled training examples (seed 42). Labels apply to the complete reviews; excerpts are shortened.}\label{tab:examples}
\begin{tabularx}{\linewidth}{@{}lrlY@{}}\toprule
Review ID & Words & Label & Excerpt\\\midrule
\code{train_09218} & 333 & Negative &
``This Hal Roach comedy short, A Tough Winter, is the ninety-ninth in the `Our Gang/Little Rascals' series and the eleventh talkie.''\\
\code{train_17779} & 316 & Positive &
``This is a very good Spanish movie but I am worried that many out there will feel lost''\ldots\\\bottomrule
\end{tabularx}
\end{table}
The examples include narrative description and mixed judgments, illustrating
why short fragments cannot reliably determine the complete review's label.
Exploratory examples were sampled from training rather than selected from test.

\subsection{Exported representations and verification}
Each JSONL record contains \code{review_id}, source and project split fields,
\code{label}, original \code{text}, minimally cleaned \code{text_clean},
\code{tokens_classical}, joined \code{text_classical}, and token-count diagnostics.
All 49,781 records were read back and compared with the prepared in-memory data.
The checks preserve IDs, labels, order and split membership, with no empty
processed reviews. We also verify the three JSONL checksums against the
exported metadata.

No exact original, cleaned or classical text is shared across final splits.
Normalization creates two extra cleaned-text and three extra classical-text
duplicates within training, which are retained because their original texts
differ. The unchanged official test set contains 199 extra original-text,
201 cleaned-text and 203 classical-text duplicate rows. None has conflicting
labels. Such rows contribute repeatedly to the reported test metrics; exact
matching also cannot rule out near-duplicates. The resulting benchmark is thus
free of detected exact cross-split overlap, rather than proven free of every
form of related content.
\FloatBarrier
